% 2026-11-13：与既有逐篇结果配套的因果讲解；原文图号保持作者编号。
\section{先在空间里看见能量路径}
\label{sec:heating_geometry}
下面的图是本报告绘制的\textbf{物理示意，不是实验数据或计算结果}。用$N$表示入流法向、$L$表示出流、$M$表示面外电流/导场方向，以免混淆不同论文的坐标：MRX中$(N,L,M)=(R,Z,Y)$；Yoon系列中为$(x,y,z)$；Drake2009中为$(y,x,z)$。MAST两环沿$z$合并，出流主要沿$r$，导场为环向。局部几何相似并不保证边界、碰撞率与加热机制相同。
\begin{figure}[htbp]
\centering
\begin{tikzpicture}[>=Stealth,font=\small]
\fill[blue!5] (-6,-2.3) rectangle (6,2.3);
\draw[ink,thick] (-5.8,1.9) .. controls (-2,1.8) and (-1,0.1) .. (0,0)
 .. controls (1,-0.1) and (2,-1.8) .. (5.8,-1.9);
\draw[ink,thick] (-5.8,-1.9) .. controls (-2,-1.8) and (-1,-0.1) .. (0,0)
 .. controls (1,0.1) and (2,1.8) .. (5.8,1.9);
\draw[->,thick] (0,2.2)--(0,0.7);
\draw[->,thick] (0,-2.2)--(0,-0.7);
\draw[->,thick] (0.7,0)--(3.4,0);
\draw[->,thick] (-0.7,0)--(-3.4,0);
\fill[orange!25] (0,0) ellipse (0.6 and 0.35);
\node at (0,0) {X};
\node[align=center] at (-4.6,0) {下游离子\textbf{温升}\\停滞、混合、耗散};
\node[align=center] at (4.6,0) {下游离子\textbf{温升}\\需要另查分布};
\node at (1.5,2.0) {入流 $N$};
\node at (2.1,-0.45) {出流 $L$};
\node[align=center,fill=white] at (-3.4,1.55) {分离线附近：电势、\\强梯度、物种解耦};
\node[align=center,fill=white] at (3.5,-1.65) {电子热可沿场输运；\\热化可滞后于重联};
\node at (0,-2.7) {$M$垂直纸面；局地电功峰、温度峰、熵产峰无需重合};
\end{tikzpicture}
\caption{将Yoo2013/F679、Tanabe2015MAST和电子相混研究放入共同几何的阅读示意。橙区只标X点邻域，不代表所有装置的电子热点宽度；强导场MAST、零导场MRX和强导场KREHM的具体图像须分别核对。}
\end{figure}

先看重联前后磁张力如何产生出流，再看电场如何给各物种做功，最后检查流能转为压力、热流把能量带往哪里，以及何种过程消去细速度结构。这三层各有证据：场/流图约束第一层，谱线或分布的中心矩约束第二层，碰撞算子、熵预算或分布趋向Maxwell的证据约束第三层。只有温度图，不能跳过中间两步。

\subsection{双温Hall的热点：磁张力释放以后，热在哪里留下来}
Browning等（2016，\texttt{L223}）的两环合并中，环向电流形成极向重联场；重新连接的场线张力推动径向出流。离子出流在下游减速、剪切，指定的离子黏性把有序流能变成内能，因而在X点两侧形成离子双峰。电子则通过电阻和超电阻通道增热，沿场热输运又改变其温度分布。图2先读磁通和电流层怎样挤压、出流、形成小岛，再读图3的温度随时间变化，才能明白两种峰为什么不同。\citep{L223}

这里有一个容易被概括丢掉的细节：\textbf{该文电子温度由超电阻耗散强烈控制，不能简单称为“电流最大处的欧姆热点”。}超电阻与电流梯度有关，作者明确说温度峰与电流密度不必同位；中心有峰或空心可能与小岛形成、抛出有关。这是指定耗散模型的机制解释。MAST实测X点窄峰是另一项证据，二者在空间上具有联系，并不构成对超电阻系数的独立测量。图3仍是初步双温结果，不能从其峰高倒推出已验证的局地$Q_i/Q_e$。
\loc{L223}{PDF第4--7页图2、3及第5--6页双温结果；第3页式(1)--(5)先给总压力模型，不能冒充完整打印的双温扩展}

\subsection{碰撞激波的因果顺序：离子先热，电子把热送到上游}
Longcope与Bradshaw（2010，\texttt{L164}）研究重联后磁通管的沿场对撞/减速。图4右图的流动从右向左，沿途依次是电子传导前沿CF、离子子激波、离子冷却区IC。离子在子激波由超声速降至亚声速，黏性应力做功使$T_i$猛升；电子因质量小、携带的流动能少，同步减速却不立即取得同样热量。下游离子通过较慢的物种碰撞交换把能量给电子，电子高平行热导率又将热传回子激波前方，形成预热区。电子压强梯度对应的极化电场还会提前减速离子，故“前沿”并非一块与流动无关的热斑。\citep{L164}

图4示例$M_1=3.5$、入流马赫数$M_i=2.4$，两者不是重联率或Alfvén马赫数。离子Prandtl数$0.33$，电子$Pr_e^{-1}=77$，解释了窄离子子激波与长电子前沿；图5比较三段空间尺度，图9--11说明它们怎样逐渐形成。因此同一位置先看到$T_i>T_e$、晚时趋同，既可能是能量交换和输运，也不能被误解为Hall项在不同时刻“换了加热对象”。该模型用共同体流、经典碰撞输运，不能自行检验MRX无碰撞层的轨道混合；文中也检验热流和离子黏性应力的适用性，不能仅因有激波解就默认流体闭合在任意薄层都成立。
\loc{L164}{PDF第10--12页式(4)--(18)，第16--20页图4、5及物理解释；后续图6、7检查热流/应力，图9--11给非稳态过程}

\subsection{电场梯度怎样破坏离子的漂移近似}
Yoon与Bellan（2018，\texttt{Yoon2018Heating}；2021，\texttt{L272}）强调的是\emph{电场空间变化相对回旋响应有多快}，不是只比较电场幅值。给定近均匀$\B=B_0\hat z$，低频垂直电场的首阶响应为$\vct v_E=\E\times\B/B^2$，下一阶极化漂移为$\vct v_p=m_i\dd\E_\perp/(q_iB^2\dd t)$。梯度很强时，沿轨道看到的场变化使这个“下一阶”不再小，离子无法保持稳定回旋中心和磁矩。局地量级指标为
\begin{equation}
\alpha_i=\frac{m_i|\nabla_\perp\cdot\E_\perp|}{|q_i|B^2},\qquad
\Omega_{ci}=\frac{|q_i|B}{m_i}.
\end{equation}
它无量纲；增大质量电荷比使失磁化更容易，增大局地$B$倾向于恢复磁化。\textbf{但绝对值指标是特定几何的诊断，不是对任意电场保证正熵产的充分判据。}例如只取一维局部$E_x=Gx$，相邻轨道满足
\begin{equation}
\delta\dddot x=\left(\frac{q_iG}{m_i}-\Omega_{ci}^2\right)\delta\dot x.
\end{equation}
此式中梯度的\emph{符号}决定局地指数不稳定性，不能用$|G|$抹掉恢复力与发散力的区别；无界轨道指数分离也不是有界混沌或热化的完整证明。这是本报告从运动方程作的局部教学推导，须与原文数值轨道证据分开。\citep{Yoon2018Heating,L272}

2018年的电子流体场呈入流向势谷、出流向势丘；离子跌入势谷后沿出流势降增能，增长场使初始相邻离子的经历不同。图1是势结构，图2是轨道，图3是流出速度展宽，三者不能用一张温度图替代。固定离子EMHD只生成背景场，测试离子无反馈，故不能同时闭合离子与电子能量。2019年\texttt{F743}的PIC图1--5进一步对照$\alpha$、温度和粒子轨道，并扫描$B_g/B_0=0$至0.5；较强导场抑制该机制，但残余平行温升区也提示存在其他通道。该趋势不能外推为所有强导场装置的离子温度都必须降低。\citep{F743}
\loc{Yoon2018Heating / L272 / F743}{依次图1--3及式(2)--(17)；正式版PDF第3页式(1)--(13)、图1；2019图1--5及静态Harris阈值的适用性讨论}

\subsection{补上的拾取参照：Drake2009与Yoon2023并非完全同一解释}
Drake等（2009，\texttt{Drake2009Pickup}）的反平行PIC显示：离子跨过狭窄出流边界时，磁场被Hall面外场迅速转向，粒子尚未来得及跟随$\E\times\B$速度。进入快速运动的磁通管后，它相对于出流具有大速度，表现为拾取状回旋/反向束。图2把轨道、磁矩破坏和分布放在一起，图3变换到随出流运动的de Hoffmann--Teller参考系，使局地电场近零；在该系中能量近守恒而运动方向改变，变回X线系后得到体流加速和中心矩展宽。\citep{Drake2009Pickup}

冷入流、近稳态、弱导场近似给出原文式(7)、(8)：
\begin{equation}
\Delta T_\parallel\simeq m_i v_0^2\frac{B_{0L}^2}{B_{0L}^2+B_{0M}^2},\quad
\Delta T_\perp\simeq\frac{m_i v_0^2}{2}\frac{B_{0M}^2}{B_{0L}^2+B_{0M}^2},\quad
\Delta T_i\simeq\frac{m_i v_0^2}{3}.
\end{equation}
温度以能量计，$v_0$为出流速度，$B_{0M}$是下游Hall场，不是可任意替换的外加导场。这是二阶矩预测，不要求Maxwell分布；文中PIC质量比25，且预测的各向异性并未精确吻合。附录A的“反弹频率”定义为$v_0/\Delta$，仅在指定极限才是实际频率，不能当成一般精确周期。

\textbf{必须保留的研究分歧。}Drake2009第16段在下游一维近稳态模型中认为，面内静电场在往返轨道上增能后又回收，面外重联电场对净位移做功。Yoon等（2023，\texttt{Yoon2023PhaseSpace}）第5节则批评把拾取作为零/弱导场重联的总解释，认为其模拟只有初始短暂拾取状阶段，随后主要是单阱中的轨道混合；作者还在结论中将净增能归于面内场，称重联电场净功为零。\emph{这是原文之间的归因差异，不能被报告写成已统一的事实。}2023图6的四阶段与图7附录的分力速度积分是具体证据，但速度分量末值恢复并不单独推出该电场的累计功为零。可检验的共同诊断应是对同一粒子、同一时间窗计算$\int q_i E_Nv_N\dd t$、$\int q_i E_Lv_L\dd t$、$\int q_i E_Mv_M\dd t$，并核对总动能增量；磁力只改变方向，分量间能量转移相消。本报告未复跑两文轨道，不能裁决其全部差异。\citep{Yoon2023PhaseSpace}
\loc{Drake2009Pickup / Yoon2023PhaseSpace}{前者全文14页含附录A；图1--4、式(7)--(10)、第16及31--39段；后者PDF第6--8页图6、7及第5节，亦见图3的冷多束}

\subsection{电子为何在重联以后、沿分离线才热化}
Zocco与Schekochihin（2011，\texttt{E512}）及Loureiro等（2013，\texttt{Loureiro2013Heating}）的强导场图像是：重联先建立流动和扰动，沿磁力线速度不同的电子逐渐错相；同一空间位置的分布在$v_\parallel$方向出现越来越细的纹理。用Hermite阶数$m$表示这些速度结构时，能量从低阶的流速、温度矩流向高阶；有限碰撞对细速度结构作用更强，才将这份非平衡自由能转为背景热。\citep{E512,Loureiro2013Heating}

以原文自由能的体积平均约定记
\begin{equation}
H_e=\frac1V\int\dd^3x\dd^3v\frac{T_{0e}g_e^2}{2F_{0e}},\qquad
\frac{\dd W_{\rm fluid}}{\dd t}=-\mathcal Q-D_\eta,\qquad
\frac{\dd H_e}{\dd t}=\mathcal Q-D_{\rm coll}.
\end{equation}
这里$F_{0e}$是背景Maxwell分布，$g_e$是除去低阶流体响应后的扰动；$\mathcal Q$是压力功导致的可逆交换（可变号），$D_\eta$是电阻耗散，$D_{\rm coll}=-V^{-1}\int T_{0e}g_e C[g_e]/F_{0e}\dd^3x\dd^3v\geq0$由满足H定理的碰撞算子给出。严格零碰撞时，$H_e$增加仍不意味着精细熵增加。\textbf{这里写的是扰动自由能预算，不能与第2章完整物种内能重复相加。}

Loureiro2013图2左列叠加磁通等值线与电流，中列叠加电势与电子平行温度扰动，右列为耗散；上、下两行分别在重联峰$22.7\tau_A$与耗散峰$29.4\tau_A$，应横向、纵向都比较。图4底行显示速度细纹形成，上两行显示自由能到高$m$再耗散，因而“先重联、后热化”具有可见证据。Tassi等（2018，\texttt{L244}）有限矩热流把这条链压缩为即时闭合，图4却缺少同样的峰时滞后，原文将是否过耗散列为问题。Numata与Loureiro（2015，\texttt{Numata2015Heating}）恢复双物种GK后，高$\beta$下离子平行/垂直相混更重要；电子多在分离线/岛缘、离子可在次级磁岛内部耗散。因此“电子用Landau、离子用Hall”不是完整分类：空间结构、相空间通道和排序共同决定可见结果。\citep{L244,Numata2015Heating}
\loc{E512 / Loureiro2013Heating / L244 / Numata2015Heating}{依次PDF第11--12页式(76)--(84)；图1、2、4、5；图4及结论；式(2.5)、图5--7与11--12、附录A}
